\DocumentMetadata{lang=es-DO,testphase={phase-III,math}}
\documentclass[11pt]{scrartcl}
\usepackage{icv}

% -----------------------------------------------------------------------
% Migrado de legacy/Acueductos y Alcantarillados/Handouts/01_metodo_aritmetico.tex
% -- primer documento real convertido al sistema de estilo icv. Ajusta
% periodo/version al ciclo real cuando lo asignes a una clase.
% -----------------------------------------------------------------------
\icvsetup{
  asignatura  = {ICV 442-T -- Acueductos y Alcantarillados},
  titulocorto = {Método aritmético},
  titulo      = {Método del incremento aritmético},
  subtitulo   = {Proyección de población},
  profesor    = {Dr. Ricardo Hernández Moreira},
  periodo     = \icvciclo{2026}{3},
  version     = {v1},
}

\begin{document}
\icvmaketitle

\begin{icvobjectives}
  \item Aplicar el método del incremento aritmético para proyectar
    población futura a partir de una serie censal.
  \item Calcular el incremento absoluto medio, con y sin intervalos
    censales regulares.
\end{icvobjectives}

\begin{icvnote}[Propósito]
Proyectar la población futura suponiendo que el incremento absoluto de
población es constante.
\end{icvnote}

\begin{icvnote}[Cuándo conviene]
Ciudades grandes, consolidadas y con crecimiento relativamente estable.
\end{icvnote}

\section{Datos ejemplo}

Se usará una serie censal con intervalos de \num{10} años. En este
documento, el incremento absoluto medio se expresará en habitantes por año.

\begin{center}
\begin{tabular}{l S[table-format=6.0] S[table-format=6.0] S[table-format=6.0]
                  S[table-format=6.0] S[table-format=6.0] S[table-format=6.0]}
  \toprule
  Año & {1970} & {1980} & {1990} & {2000} & {2010} & {2020} \\
  \midrule
  Población & 112400 & 126100 & 139500 & 153400 & 167000 & 180600 \\
  \bottomrule
\end{tabular}
\end{center}

\section{Ecuaciones mínimas}

El modelo diferencial es:
\begin{equation}
  \frac{dP}{dt}=K
\end{equation}
donde \(P\) es la población y \(K\) es el incremento absoluto constante,
expresado aquí en habitantes por año.

La ecuación de proyección es:
\begin{equation}
  P(t)=P_0+Kt
\end{equation}
donde \(P_0\) es la población al inicio del intervalo de proyección y
\(t\) es el tiempo transcurrido en años.

Si los intervalos censales son regulares, el incremento medio anual se
calcula como:
\begin{equation}
  K=\frac{1}{N}\sum_{i=1}^{N} k_i
\end{equation}
donde \(k_i=\Delta P_i/\Delta t_i\) es el incremento anual del período
censal \(i\), \(\Delta P_i\) es el incremento de población en ese período,
\(\Delta t_i\) es su duración en años y \(N\) es el número de períodos
usados.

\begin{icvwarning}[Nota sobre intervalos censales irregulares]
Si los intervalos censales no son regulares, conviene usar el promedio
ponderado
\[
  K_w=\frac{\sum_{i=1}^{N} \Delta t_i\,k_i}{\sum_{i=1}^{N}\Delta t_i}
\]
con pesos iguales a la duración de cada período en años.
\end{icvwarning}

\section{Procedimiento}
\begin{enumerate}
  \item Calcular el incremento absoluto de población en cada período censal.
  \item Dividir cada incremento entre la duración del período para obtener
    incrementos anuales.
  \item Promediar esos incrementos anuales.
  \item Sustituir el valor de \(K\) en la ecuación de proyección.
\end{enumerate}

\section{Ejemplo resuelto}

\begin{center}
\begin{tabular}{l S[table-format=6.0] S[table-format=5.0] S[table-format=4.0]}
  \toprule
  Año & {Población} & {Incremento del período} & {Incremento anual} \\
  \midrule
  1970 & 112400 & {--} & {--}  \\
  1980 & 126100 & 13700 & 1370 \\
  1990 & 139500 & 13400 & 1340 \\
  2000 & 153400 & 13900 & 1390 \\
  2010 & 167000 & 13600 & 1360 \\
  2020 & 180600 & 13600 & 1360 \\
  \midrule
  \multicolumn{3}{l}{Promedio} & 1364 \\
  \bottomrule
\end{tabular}
\end{center}

Entonces,
\[
  K=\frac{\num{1370}+\num{1340}+\num{1390}+\num{1360}+\num{1360}}{5}
    =\num{1364}\ \text{hab/año}
\]

Como 2030, 2040 y 2050 están a 10, 20 y 30 años de 2020, respectivamente:
\[
  P_{2030}=\num{180600}+\num{1364}(10)=\num{194240}
\]
\[
  P_{2040}=\num{180600}+\num{1364}(20)=\num{207880}
\]
\[
  P_{2050}=\num{180600}+\num{1364}(30)=\num{221520}
\]

\begin{icvwarning}[Resultado clave]
La población proyectada es \num{194240} habitantes en 2030,
\num{207880} habitantes en 2040 y \num{221520} habitantes en 2050.
\end{icvwarning}

\section{Teoría breve}

Partiendo de la ecuación diferencial:
\[
  \frac{dP}{dt}=K
\]
se separan variables:
\[
  dP=K\,dt
\]
y se integra:
\[
  \int dP=\int K\,dt
\]
de donde resulta:
\[
  P=Kt+C_1
\]
Si en \(t=0\) la población es \(P_0\), entonces \(C_1=P_0\) y se obtiene
la ecuación de proyección:
\[
  P(t)=P_0+Kt
\]

\section{Observaciones de uso}
\begin{itemize}
  \item Tiende a subestimar ciudades jóvenes o de expansión acelerada.
  \item Es útil como estimación conservadora.
  \item Debe indicarse siempre si \(K\) está expresado por año, por
    década o por otro período censal.
\end{itemize}

\end{document}
